\chapter{Die Zeit, die falsch gewählt war}\label{fall:zeit}

\begin{merksatz}
\boxtitel{Die Vorlage}
\textbf{Zusammenhang.} Naht zwischen Speicherung und Auswertung. Kein Code --- eine Feldbelegung, die seit acht Jahren gilt.

\begin{center}
\begin{minipage}{0.9\textwidth}
\code{Befund.zeitpunkt}\ \ \ \ \ \code{-- gefuellt mit: Zeitpunkt}\\
\code{}\ \ \ \ \ \ \ \ \ \ \ \ \ \ \ \ \ \ \ \ \ \ \ \code{\ \ \ des Datenbankschreibens}
\end{minipage}
\end{center}

\begin{center}
\renewcommand{\arraystretch}{1.45}
\begin{tabular}{@{}P{3.0cm}P{2.2cm}P{2.6cm}@{}}
\toprule
\textbf{Eigenschaft} & \textbf{Bewertung} & \\
\midrule
Immer vorhanden & \textcolor{gruen}{\textbf{ja}} & \\
Immer aufsteigend & \textcolor{gruen}{ja} & \textcolor{gruen}{sortierbar}\\
Kostenlos & \textcolor{gruen}{ja} & \textbf{Datenbank setzt ihn}\\
Nie leer & \textcolor{gruen}{ja} & \\
\bottomrule
\end{tabular}
\end{center}

\begin{center}
\begin{minipage}{0.88\textwidth}
\kursiv{Vier gute Eigenschaften --- und alle vier sind technisch. Der Schreibzeitpunkt ist der zuverlässigste Wert im ganzen System und beantwortet die Frage nicht, die man ihm stellt.}
\end{minipage}
\end{center}
\end{merksatz}

\begin{merksatz}
\boxtitel{Die Diagnose}
\textbf{Abstand zwischen Abnahme und Speicherung.}

\begin{center}
{\small
\renewcommand{\arraystretch}{1.45}
\begin{tabular}{@{}P{2.8cm}P{2.0cm}P{2.0cm}P{2.0cm}@{}}
\toprule
\textbf{Herkunft} & \textbf{Median} & \textbf{$95.$ Perz.} & \textbf{Anteil}\\
\midrule
Bettseitig & $\mathbf{40}$ s & $3$ min & $31\,\%$\\
Zentrallabor & $52$ min & $2{,}1$ h & $54\,\%$\\
Mikrobiologie & $\mathbf{2{,}3}$ \textbf{d} & $6$ d & $9\,\%$\\
\textbf{Fremdhaus} & \textcolor{akzent}{$\mathbf{4}$ \textbf{d}} & \textcolor{akzent}{$\mathbf{31}$ \textbf{d}} & $6\,\%$\\
\bottomrule
\end{tabular}}
\end{center}

\begin{center}
\begin{minipage}{0.88\textwidth}
\kursiv{Bei einem Drittel der Befunde beträgt der Abstand unter drei Minuten --- dort ist der Schreibzeitpunkt eine brauchbare Näherung. Bei fünfzehn Prozent liegt er bei Tagen, und dort ist der Verlauf falsch.}
\end{minipage}
\end{center}

\textbf{Wie sich das auswirkt.}

\begin{center}
\renewcommand{\arraystretch}{1.45}
\begin{tabular}{@{}P{3.0cm}P{2.2cm}P{2.6cm}@{}}
\toprule
\textbf{Auswirkung} & \textbf{Häufigkeit} & \\
\midrule
Verlauf falsch sortiert & $\approx\mathbf{4\,\%}$ & \textbf{der Ansichten}\\
Steigungs\-warnung falsch & $\approx\mathbf{1{,}2\,\%}$ & \textcolor{akzent}{\textbf{je Nacherfassung}}\\
Verlauf am Stichtag & \textcolor{akzent}{nicht rekonstruierbar} & \textcolor{akzent}{\textbf{regulatorisch}}\\
\bottomrule
\end{tabular}
\end{center}

\begin{center}
\kursiv{Die dritte Zeile wiegt am schwersten und fiel erst bei einer Prüfung auf: Die Frage \enquote{welche Werte lagen der Ärztin am Dienstag vor} lässt sich nicht beantworten, weil es nur eine Zeitachse gibt --- Band 23, Kapitel~4.}
\end{center}
\end{merksatz}

\begin{technik}
\boxtitel{Das Urteil: überarbeitet}
\begin{center}
\begin{minipage}{0.9\textwidth}
\code{Befund.abnahme\ \ \ \ \ \ \ Pflicht\ \ fachliche Zeit}\\
\code{Befund.eingang\ \ \ \ \ \ \ Pflicht}\\
\code{Befund.validiert\_am\ \ optional}\\
\code{Befund.geschrieben\ \ \ Pflicht\ \ Systemzeit, unveraenderlich}\\
\code{Befund.zeit\_unsicher bool\ \ \ \ \ Abnahme geschaetzt?}
\end{minipage}
\end{center}

\begin{center}
\begin{minipage}{0.86\textwidth}
\kursiv{Vier Zeitpunkte statt einem, und ein Kennzeichen. Das letzte Feld ist die ehrliche Antwort auf den Altbestand: Wo die Abnahmezeit nicht bekannt ist, wird sie nicht geschätzt, sondern als unsicher gekennzeichnet.}
\end{minipage}
\end{center}
\end{technik}

\begin{merksatz}
\boxtitel{Die Begründung}
\begin{center}
{\small
\renewcommand{\arraystretch}{1.45}
\begin{tabular}{@{}P{2.8cm}P{2.0cm}P{2.2cm}P{2.0cm}@{}}
\toprule
& \textbf{Vorlage} & \textbf{Überarbeitet} & \\
\midrule
Zeitfelder & $1$ & $\mathbf{4}$ & \\
Byte je Befund & $8$ & $\mathbf{33}$ & \\
Bei $10^9$ Befunden & --- & \textcolor{mittel}{$+25$ GB} & \textcolor{gruen}{tragbar}\\
Verlauf richtig & $\approx96\,\%$ & \textcolor{gruen}{$\approx\mathbf{100\,\%}$} & \\
Stichtag rekonstruierbar & \textcolor{akzent}{\textbf{nein}} & \textcolor{gruen}{\textbf{ja}} & \\
Umstellung & --- & \textcolor{akzent}{$\mathbf{9}$ \textbf{Wochen}} & \\
Aufwand & --- & $\mathbf{22}$ PT & \\
\bottomrule
\end{tabular}}
\end{center}

\begin{center}
\begin{minipage}{0.88\textwidth}
\kursiv{Zweiundzwanzig Personentage und neun Wochen Umstellung für vier Zeitfelder --- der teuerste Fall dieses Bandes. Er ist es wert, weil die vorletzte Zeile eine regulatorische Pflicht betrifft, die sich nachträglich nicht herstellen lässt.}
\end{minipage}
\end{center}

\textbf{Der Altbestand.}

\begin{center}
\renewcommand{\arraystretch}{1.45}
\begin{tabular}{@{}P{3.0cm}P{2.2cm}P{2.6cm}@{}}
\toprule
\textbf{Umgang} & \textbf{Gewählt?} & \\
\midrule
Abnahme aus Schreibzeit ableiten & \textcolor{akzent}{\textbf{nein}} & \textcolor{akzent}{\textbf{erfundene Daten}}\\
Alte Werte löschen & \textcolor{akzent}{nein} & \\
\textbf{Als unsicher kennzeichnen} & \textcolor{gruen}{\textbf{ja}} & \textcolor{gruen}{$412$ Mio.\ Werte}\\
Stichtag nennen & \textcolor{gruen}{ja} & \textbf{ab da verlässlich}\\
\bottomrule
\end{tabular}
\end{center}

\begin{center}
\kursiv{Die erste Zeile war der naheliegende Vorschlag und wurde verworfen: Aus dem Schreibzeitpunkt eine Abnahmezeit zu berechnen erzeugt Werte, die richtig aussehen und es nicht sind. Vierhundertzwölf Millionen Befunde tragen jetzt ein Kennzeichen --- unbequem und ehrlich.}
\end{center}
\end{merksatz}

\begin{warnung}
\boxtitelr{Der Fehlschluss}
\begin{center}
\renewcommand{\arraystretch}{1.5}
\begin{tabular}{@{}P{3.2cm}P{2.2cm}P{2.4cm}@{}}
\toprule
\textbf{Naheliegend} & \textbf{Richtig?} & \\
\midrule
\enquote{Immer alle Zeitpunkte speichern} & \textcolor{akzent}{\textbf{nein}} & \textcolor{akzent}{was gebraucht wird}\\
\enquote{Der Schreibzeitpunkt ist wertlos} & \textcolor{akzent}{\textbf{nein}} & \textcolor{gruen}{\textbf{er bleibt}}\\
\enquote{$96\,\%$ richtig ist gut genug} & \textcolor{akzent}{nein} & \textbf{die $4\,\%$ sind sortiert}\\
\textbf{\enquote{Ein Feld ist keine Bedeutung}} & \textcolor{gruen}{\textbf{ja}} & \\
\bottomrule
\end{tabular}
\end{center}

\begin{center}
\begin{minipage}{0.88\textwidth}
\kursiv{Die zweite Zeile ist wichtig: Der Schreibzeitpunkt wird nicht ersetzt, sondern ergänzt --- er ist die Systemzeit und beantwortet die Nachvollzugsfrage. Falsch war nicht das Feld, sondern dass es als fachliche Zeit gelesen wurde.}
\end{minipage}
\end{center}

\begin{center}
\kursiv{Die dritte Zeile ist die statistische Falle: Die vier Prozent falsch sortierter Verläufe sind nicht zufällig verteilt --- sie häufen sich bei Mikrobiologie und Fremdbefunden, also genau dort, wo eine Nacherfassung stattfindet und ein Verlauf klinisch etwas bedeutet.}
\end{center}
\end{warnung}
\vspace{-4pt}
\noindent Klepp\-mann\zit{18} unterscheidet die beiden Zeit\-ach\-sen und ihre Ver\-wen\-dung; Evans\zit{9} begründet, warum ein Feld\-na\-me eine fach\-li\-che Aus\-sa\-ge trägt.
